\newtoggle{withbonus} \toggletrue{withbonus} % set true to show bonus points in grade table

% setting underlying course name (Grundlagenfach Mathematik, §-Unterricht, \ldots)
\newcommand{\mycourse}{Ergänzungsfach Informatik}
% name of the class
\newcommand{\myclass}{4. Klassen}
% set date
\newcommand{\mydate}{20. November 2025}

\title{\textbf{Kolmogorov-Komplexität}}
\author{Thomas Graf\\
	\mycourse, \myclass}
\date{\mydate}
\maketitle
\thispagestyle{headandfoot}

\begin{center}
	\fbox{\parbox{140mm}{\centering
			\begin{itemize}
				\item Dauer der Prüfung: $65$ Minuten
				      %\item \textbf{Zeige in jeder Aufgabe unbedingt Deine Schritte!}
				\item Erlaubte Hilfsmittel:
				      \begin{itemize}
					      \item Schreibzeug
					            % \item Taschenrechner
				      \end{itemize}
			\end{itemize}}}
\end{center}
\vspace{10mm}

\begin{center}
	Vorname: \rule{45mm}{0.8pt}\qquad Nachname: \rule{45mm}{0.8pt}
\end{center}
\vspace{20mm}

\begin{center}
	\iftoggle{withbonus} {
		\combinedgradetable[h][questions]
	} {
		\gradetable[h][questions]
	}
\end{center}
\vspace{15mm}
\begin{center}
	Note:
	\vspace{10mm}

	\rule{8mm}{1.2pt}
\end{center}
% \thispagestyle{empty}
\clearpage


\begin{questions}


\question[2] Erkläre die Begriffe \texttt{WYSIWYG} und \texttt{WYSIWYM}.
	\begin{solution}
		TODO
	\end{solution}

\question[2] Erstelle ein neues \LaTeX-Macro mit dem Namen \lstinline[style=LaTeX]|coloredmath|, sodass der Befehl \lstinline[style=LaTeX]|\coloredmath{red}{\sqrt{5}}| zum Resultat
\begin{center}
	\textcolor{Red}{$\sqrt{5}$}
\end{center}
kompiliert.

\question[2]

Wir schreiben ein \LaTeX-Dokument, welches zu folgendem Text kompilieren soll:
\begin{adjustwidth}{1em}{2em}
\section*{Städtetourismus am Ende? In Europas Metropolen macht sich Verzweiflung breit (\textit{Quelle: NZZ})}
In den Herbstferien unternehmen Schweizer gerne Städtereisen. Mittlerweile haben aber der Bundesrat und andere Regierungen diverse Metropolen zu Risikogebieten erklärt.
\subsection*{Paris: Das Aufatmen im Sommer war nur kurz}
\begin{itemize}
    \item Etwas mehr als die Hälfte der zuletzt rund 38 Mio. Besucher kommt aus dem Ausland, die grösste Gruppe aus den USA, darauf folgen Briten und Deutsche.
    \item Die für Tourismus und Kongresse zuständige Behörde schätzt den Einbruch beim Freizeittourismus auf das Jahr 2020 gesehen auf \SI{67}{\percent}.
    \item Der Ausblick auf die Wintermonate gibt wenig Anlass zu Optimismus: Viele Nachbarländer haben die Hauptstadtregion zur Risikozone erklärt. 
\end{itemize}
\end{adjustwidth}
Finde mindestens vier syntaktische Fehler in unserem \LaTeX-Dokument. Für jeden gefundenen Fehler gibt es einen halben Punkt (maximal 2 Punkte):
\begin{lstlisting}[style=LaTeX]
\usepackage{siunitx}
\documentclass{article}
\section{Städtetourismus am Ende? In Europas Metropolen macht sich Verzweiflung breit (Quelle: NZZ)}
In den Herbstferien unternehmen Schweizer gerne Städtereisen. Mittlerweile haben aber der Bundesrat und andere Regierungen \textbf{diverse} Metropolen zu Risikogebieten erklärt.
\subsection*{Paris: Das Aufatmen im Sommer war nur kurz}
\begin{item}
    \item Etwas mehr als die Hälfte der zuletzt rund 38 Mio. Besucher kommt aus dem Ausland, die grösste Gruppe aus den USA, darauf folgen Briten und Deutsche.
    \item Die für Tourismus und Kongresse zuständige Behörde schätzt den Einbruch beim Freizeittourismus auf das Jahr 2020 gesehen auf \SI{67}{\percent}.
    \item Der Ausblick auf die Wintermonate gibt wenig Anlass zu Optimismus: Viele Nachbarländer haben die Hauptstadtregion zur Risikozone erklärt. 
\end{item}
\end{article}
\end{lstlisting}

\begin{solution}
Der korrekte Code, der den Text in der Aufgabenstellung erzeugt ist:
\begin{lstlisting}[style=LaTeX]
\documentclass{article}
\begin{document}
\usepackage{siunitx}
\section*{Städtetourismus am Ende? In Europas Metropolen macht sich Verzweiflung breit (\textit{Quelle: NZZ})}
In den Herbstferien unternehmen Schweizer gerne Städtereisen. Mittlerweile haben aber der Bundesrat und andere Regierungen diverse Metropolen zu Risikogebieten erklärt.
\subsection*{Paris: Das Aufatmen im Sommer war nur kurz}
\begin{itemize}
    \item Etwas mehr als die Hälfte der zuletzt rund 38 Mio. Besucher kommt aus dem Ausland, die grösste Gruppe aus den USA, darauf folgen Briten und Deutsche.
    \item Die für Tourismus und Kongresse zuständige Behörde schätzt den Einbruch beim Freizeittourismus auf das Jahr 2020 gesehen auf \SI{67}{\percent}.
    \item Der Ausblick auf die Wintermonate gibt wenig Anlass zu Optimismus: Viele Nachbarländer haben die Hauptstadtregion zur Risikozone erklärt. 
\end{itemize}
\end{document}
\end{lstlisting}
Fehler:
\begin{itemize}
	\item \lstinline[style=LaTeX]|\usepackage{siunitx}| muss nach der Dokumentklasse stehen
    \item \texttt{*} bei \lstinline[style=LaTeX]|\section| fehlt
    \item textit im Titel fehlt
    \item \lstinline[style=LaTeX]|\end{document}| statt \lstinline[style=LaTeX]|\end{article}|
    \item \lstinline[style=LaTeX]|\begin{itemize} .. \end{itemize}| statt \lstinline[style=LaTeX]|\begin{item} .. \end{item}|
    \item textbf zu viel
\end{itemize}
\end{solution}

	\question[2] Wir haben die Kolmogorov-Komplexität für binäre Wörter definiert. Wie wird die Kolmogorov-Komplexität für eine natürliche Zahl $n$ definiert?


	\question[2] Bestimme die Länge der kürzesten Darstellung der Zahl $54345120$ in Basis $6$.

	\begin{solution}
		TODO $\ceil{\log_2\lr{92340283 + 1}} = 27$.
	\end{solution}

	\question[2] Sei $w_n := 0^{2^{n^3}}$ für jedes $n\in\N, n>0$ ein binäres Wort.	Gib eine (asymptotisch) möglichst gute obere Schranke für die Kolmogorov-Komplexität $K(w_n)$ von $w_n$ an, gemessen in der Länge $\abs{w_n}$ von $w_n$.
	\begin{solution}
		Das folgende \texttt{C'++} Programm erzeugt das Wort $w_n$:
		\begin{lstlisting}[language=C++,caption=$W_n$]
    #include <iostream>
    #include <cmath> // für die pow Funktion
    
    void Wn() {
        int k = ***n+++;
        int k = pow(k, 3);
        int k = pow(2, k);
        for (int i = 1; i <= k; ++i) {
            std::cout << 0;
        }
        std::cout << std::endl;
    }
\end{lstlisting}
		TODO
	\end{solution}

	\question[2] Gibt eine Funktion $f: \N\to\N$ an, sodass für jede natürliche Zahl $n$ mindestens drei Viertel der natürlichen Zahlen kleiner als $f(n)$ eine Kolmogorov-Komplexität von mindestens $2n$ haben.

    Tipps:
    \begin{itemize}
        \item Wie viele unterschiedliche nichtleere binäre Wörter der Länge $2n$ kann es höchstens geben?
        \item Wie viele natürliche Zahlen mit einer Kolmogorov-Komplexität von weniger als $2n$ gibt es also?
        \item Wähle $f(n)$ als das mindestens Vierfache dieser Anzahl.
    \end{itemize}
	\begin{solution}
		TODO
	\end{solution}
	
	\question[2] Im Skript haben wir bewiesen:
	\begin{center}
        Das Problem, für jedes binäre Wort $w$ die Kolmogorov-Komplexität $K(w)$ zu berechnen, ist algorithmisch unlösbar (nicht berechenbar).
	\end{center}
    Erkläre die zentralen Ideen dieses Beweise in deinen eignen (präzisen) Worten.
	\begin{solution}
		Siehe das Theorem zur Nichtberechenbarkeit (im allgemeinen) der Kolmogorov-Komplexität im Skript.
	\end{solution}


	\question[2] Gib eine streng monoton wachsende Folge $y_1,y_2,\ldots$ natürlicher Zahlen an, sodass eine Konstante $c$ existiert, sodass für alle $i\geq 1$
	\begin{align*}
		K\lr{y_i} \leq \ceil{\log_2\lr{\log_2\lr{\log_2\lr{\sqrt[4]{y_i}}}}} + c
	\end{align*}
	gilt.

\end{questions}
